\documentclass[10pt,a4paper]{amsart}
\usepackage{amsmath,amssymb,amsthm,graphicx}
\usepackage{amsmath}
\usepackage{amssymb}
\usepackage{hyperref}
\newtheorem{Theorem}{Theorem}[section]
\newtheorem{Corollary}[Theorem]{Corollary}
\newtheorem{Lemma}[Theorem]{Lemma}
\newtheorem{Proposition}[Theorem]{Proposition}
{\theoremstyle{definition}
\newtheorem{Note}[Theorem]{Note}
\newtheorem{Definition}[Theorem]{Definition}
\newtheorem{Problem}[Theorem]{Problem}
}
\usepackage{cleveref}
\crefname{originalproposition}{original proposition}{original propositions}
\Crefname{originalproposition}{Original proposition}{Original propositions}
\makeatletter
\newlabel{lem:com2}{{1.54}{source}{}{sourcepointer-1}{}}
\newlabel{lem:com2@cref}{{[originalproposition][0][]1.54}{[1][0][]source}{}{}{}}
\makeatother
\title{The two-sided ideal generated by the commutators of the Askey-Wilson algebra coincides with the kernel of its polynomial quotient}
\author{Paul Terwilliger}
\date{doi:10.3842/SIGMA.2011.069; CC BY4.0}
\begin{document}
\maketitle
\noindent\textit{Editorial note.} All delivered mathematical text is the author's own, copied from the pinned publisher source sources/sigma-2011-069-askey-wilson/pinned-original-source.tex; no statement or proof step was rewritten, reordered or removed.\par\medskip
\noindent\textit{Editorial note.} Preamble provenance: the retained preamble is the source's own layout adapter plus the author's macro block at lines 10-18 of sources/sigma-2011-069-askey-wilson/pinned-original-source.tex (everything between the publisher class metadata and \textbackslash{}begin\{document\}, comments excluded). Nothing was added, no macro was invented, and the publisher document class is deliberately not reproduced; the excerpt serves only as preamble, never as a span coordinate.\par\medskip
\noindent\textit{Editorial note.} External references are navigation pointers only; every deliberately kept step of the author's argument stays inside the counted spans. The displayed numbers of those pointers are the author's own numbering under the author's \textbackslash{}newtheorem declarations (the source has no subsections, so numbering is section.index), computed mechanically by the lane's scanner.\par\medskip
\noindent\textit{Editorial note.} The counted statement itself cites Lemma lem:com2 (line 2806); that label is recorded once as a navigation stub of kind proposition, and every step that consumes the construction stays inside the counted proof.\par\medskip
\noindent\textit{Editorial note.} The counted spans contain no bibliography citation, so no bibliography excerpt is attached to this card.\par\medskip
\noindent\textit{Editorial note.} Licence binding: the card's source\_version now carries the value recorded in this source's own licence evidence (sources/sigma-2011-069-askey-wilson/license-observation.json), namely 'SIGMA7(2011),069; publisher final form2011-07-09'. The adopt step of the publication line compares the two and refused the earlier revision; the card is aligned to the evidence rather than the evidence to the card.\par\medskip
\noindent\textit{Allowed standard prerequisites (editorial).}\begin{itemize}
\item The homomorphism from the universal Askey-Wilson algebra onto the polynomial algebra, and its surjectivity. Same source, Lemma labelled lem:com2; invoked in the counted statement at line 2806 and inside the counted proof at lines 2811-2818. Supplies the map whose kernel Gamma is one of the two objects compared in the statement.
\item The monomials in the barred generators form a basis of the polynomial algebra. Used inside the counted proof at lines 2835-2847. Gives that the monomials A\textasciicircum{}i B\textasciicircum{}j C\textasciicircum{}k form a complement M of Gamma in the algebra, so that the sum Delta = M + Gamma is direct.
\end{itemize}\medskip
\ifdefined\hypertarget\hypertarget{sourcepointer-1}{}\fi
\noindent\textit{Source reference.} Original proposition 1.54: sources/sigma-2011-069-askey-wilson/pinned-original-source.tex, source lines 2806--2806. Navigation pointer for \textbackslash{}ref\{lem:com2\} used inside the counted span (line 2806); the argument that consumes it is retained in the span.\par\medskip
\begin{Proposition}
\label{lem:coincide}
The following coincide:
\begin{enumerate}\itemsep=0pt
\item[$(i)$]
The $2$-sided ideal
$\Delta \lbrack  \Delta ,\Delta \rbrack \Delta $;
\item[$(ii)$]
The kernel of the homomorphism
$\Delta \to
\mathbb F\lbrack
\overline A,
\overline B,
\overline C
\rbrack $ from Lemma~{\rm \ref{lem:com2}}.
\end{enumerate}
\end{Proposition}

\begin{proof}
Let $\Gamma$ denote the kernel of the
 homomorphism
$\Delta \to
\mathbb F\lbrack
\overline A,
\overline B,
\overline C
\rbrack $
 and note that~$\Gamma $ is a~2-sided ideal
 of~$\Delta $.
We have
$\lbrack \Delta ,\Delta \rbrack \subseteq \Gamma$
since
$
\mathbb F\lbrack
\overline A,
\overline B,
\overline C
\rbrack $ is commutative,
and
$\Delta \lbrack \Delta ,\Delta \rbrack \Delta
\subseteq \Gamma $
since
$\Gamma $ is a 2-sided ideal of $\Delta $.
The elements
$\lbrace \overline A^i \overline B^j \overline C^k |i,j,k\geq 0\rbrace $
form a  basis for the $\mathbb F$-vector space
$
\mathbb F\lbrack
\overline A,
\overline B,
\overline C
\rbrack $.
Therefore
the elements $\lbrace A^iB^jC^k |i,j,k\geq 0\rbrace$
form a basis for a complement of
$\Gamma $ in $\Delta $.
Denote this complement by~$M$, so
the sum
$\Delta =M+\Gamma $ is direct.
We now show $\Delta =M+
\Delta \lbrack \Delta ,\Delta \rbrack \Delta $.
The $\mathbb F$-algebra~$\Delta $
is generated by~$A$,~$B$,~$C$. Therefore
the $\mathbb F$-vector space~$\Delta $
is spanned by elements of the form
$x_1x_2\cdots x_n$ where $n\geq 0$ and
$x_i \in \lbrace A,B,C\rbrace$ for
$1 \leq i \leq n$.
Let
$x_1x_2\cdots x_n$ denote such
an element. Then
$\Delta \lbrack \Delta ,\Delta \rbrack \Delta $ contains
\begin{gather*}
x_1\cdots x_{i-1}x_i\cdots x_n  -
x_1\cdots x_ix_{i-1}\cdots x_n
\end{gather*}
for $2 \leq i \leq n$.
Therefore
$\Delta \lbrack \Delta ,\Delta \rbrack \Delta $ contains
\begin{gather*}
x_1x_2\cdots x_n  -  A^iB^j C^k,
\end{gather*}
where $i$, $j$, $k$ denote the number of
times $A$, $B$, $C$ appear among
$x_1, x_2,\ldots,  x_n$.
Observe $A^iB^jC^k \in M$ so
$x_1x_2\cdots x_n \in
M+
\Delta \lbrack \Delta ,\Delta \rbrack \Delta $.
Therefore
$\Delta =M+
\Delta \lbrack \Delta ,\Delta \rbrack \Delta $.
We already showed
$\Delta \lbrack \Delta ,\Delta \rbrack \Delta
\subseteq \Gamma $
and the sum
$\Delta =M+\Gamma $ is direct. By these comments
$\Delta \lbrack \Delta ,\Delta \rbrack \Delta
=\Gamma $.
\end{proof}

\begin{thebibliography}{99}

\footnotesize\itemsep=0pt
\bibitem{alpern}
Alperin R.C.,
Notes: ${\rm  {PSL}}_2(\mathbb Z)= \mathbb Z_2 \star \mathbb Z_3$,
\href{http://dx.doi.org/10.2307/2324963}{{\it Amer. Math. Monthly}} {\bf 100} (1993), 385--386.
\end{thebibliography}
\end{document}
